\chapter*{概要}
\addcontentsline{toc}{chapter}{概要}

実験 S1 では、小型株で、決算発表文の読み取りが発表後のドリフト（リターンの向き）を予測した。この実験 S2 では、同じ 1,000 社の決算発表文で、読み取りが発表後の\textbf{変動率（リスク）}を、過去の変動率を超えて予測するかを、事前登録して 1 回だけ確かめた。

\begin{keybox}[要点]
\begin{itemize}
\item \textbf{主な検定は支持された}。ChronoGPT（2015 年版）のベクトルに、変動率の驚きを正解にしたリッジを当てた点数が 1 標準偏差高い発表は、発表の前の変動率・2 日の反応の大きさなどを統制した後、2〜60 取引日の変動率が \V{Pcoef}\V{Punit} 高かった（$t$ = \V{Pt}、置換の片側 $p$ \V{Pperm}、発表 \V{Pn} 件）。2021 年以降では \V{Rcoef}\V{Runit}（$t$ = \V{Rt}）。
\item \textbf{ただし大きさは小さい}。決定係数の増分は \V{Prtwo} ポイントで、変動率の大部分は発表の前の変動率で説明される。
\item \textbf{ほかの点数}：ChronoGPT の 26 層目、gpt-4o-mini の 14 項目＋リッジ、辞書の否定的な語も有意に正。辞書の不確実さの語、FinBERT の否定的な文は有意でなかった。2 日の反応の大きさは、辞書の否定的な語や gpt-4o-mini のリスクの印の数がよく予測した。
\item \textbf{標本外の予測には役立たなかった}。四半期ごとに学習した単純な予測のモデルに点数を足すと、二乗誤差はわずかに大きくなった（比 \V{Fratio}、$t$ = \V{Ft}）。点数の情報は、同じ時期の銘柄どうしの相対的なリスクについてのもので、変動率の水準の予測には効かなかった。
\item \textbf{次}：BLM の $\Omega$ や量の調整に、リスクの点数を相対的に使う（新しい事前登録の追記として）。
\end{itemize}
\end{keybox}
